On the square lattice the injective regions are built from $2\times3$ and
$3\times2$ coordinate rectangles. This section fixes the geometry: the nested
regions $R\subseteq S$ with $S\setminus R$ a single site, which realize the
one-site-different injective regions the normal theorem requires at each vertex,
together with the local window $T$ used around an edge. The union closure of
Section~\ref{sec:peps_normal_union} builds these regions from injective
rectangles, so a rectangular injectivity hypothesis yields their injectivity.

\begin{definition}[Complement of a finite region]%
    \label{def:peps_regionComplement}
    \lean{TNLean.PEPS.regionComplement}
    \lean{TNLean.PEPS.mem_regionComplement}
    \leanok
    For a finite vertex region $R\subseteq V$, its complement is
    $V\setminus R$. Equivalently, a vertex lies in the complement exactly
    when it does not lie in $R$.
\end{definition}

\begin{definition}[Square-lattice coordinate regions]%
    \label{def:peps_squareLatticeCoordinateRegions}
    \lean{TNLean.PEPS.SquareLatticeVertex}
    \lean{TNLean.PEPS.squareLatticeHorizontalNeighbor}
    \lean{TNLean.PEPS.squareLatticeVerticalNeighbor}
    \lean{TNLean.PEPS.squareLatticeGraph}
    \lean{TNLean.PEPS.IsHorizontalSquareLatticeEdge}
    \lean{TNLean.PEPS.IsVerticalSquareLatticeEdge}
    \lean{TNLean.PEPS.squareLatticeRectangle}
    \lean{TNLean.PEPS.squareLatticeFull}
    \lean{TNLean.PEPS.squareLatticeCoordinateInterval}
    \lean{TNLean.PEPS.squareLatticeContiguousRectangle}
    \lean{TNLean.PEPS.IsTwoByThreeSquareLatticeProduct}
    \lean{TNLean.PEPS.IsThreeByTwoSquareLatticeProduct}
    \lean{TNLean.PEPS.IsTwoByThreeContiguousSquareLatticeRectangle}
    \lean{TNLean.PEPS.IsThreeByTwoContiguousSquareLatticeRectangle}
    \lean{TNLean.PEPS.normalSquareRegionR}
    \lean{TNLean.PEPS.normalSquareRegionS}
    \lean{TNLean.PEPS.normalSquareRegionTVerticalBlock}
    \lean{TNLean.PEPS.normalSquareRegionTHorizontalBlock}
    \lean{TNLean.PEPS.normalSquareRegionTHole}
    \lean{TNLean.PEPS.normalSquareRegionT}
    \lean{TNLean.PEPS.normalSquareEdgeComplementRegion}
    \lean{TNLean.PEPS.squareLatticeGraph_adj}
    \lean{TNLean.PEPS.squareLatticeGraph_adj_right}
    \lean{TNLean.PEPS.squareLatticeGraph_adj_up}
    \lean{TNLean.PEPS.squareLatticeRightEdge}
    \lean{TNLean.PEPS.squareLatticeUpEdge}
    \lean{TNLean.PEPS.squareLatticeRightEdge_isHorizontal}
    \lean{TNLean.PEPS.squareLatticeUpEdge_isVertical}
    \lean{TNLean.PEPS.squareLatticeEdge_horizontal_or_vertical}
    \lean{TNLean.PEPS.squareLatticeEdge_not_horizontal_and_vertical}
    \lean{TNLean.PEPS.horizontalSquareLatticeEdge_coords}
    \lean{TNLean.PEPS.horizontalSquareLatticeEdge_eq_rightEdge}
    \lean{TNLean.PEPS.verticalSquareLatticeEdge_coords}
    \lean{TNLean.PEPS.verticalSquareLatticeEdge_eq_upEdge}
    \leanok
    The finite $n\times m$ square lattice is represented by pairs of
    coordinates. The corresponding square-lattice graph has horizontal and
    vertical nearest-neighbor edges; the two elementary adjacency lemmas record
    the right-neighbor and upper-neighbor cases used for translated edge
    blockings, and the corresponding coordinate edges name those two
    nearest-neighbor edges as elements of the edge set. These named right and
    upward edges have the expected horizontal and vertical orientations.
    Every edge of this
    graph is horizontal or vertical, but not both. Since edges are represented
    by ordered endpoint pairs, a horizontal edge is oriented from left to right
    and is exactly the coordinate right edge from its left endpoint; similarly,
    a vertical edge is oriented from bottom to top and is exactly the
    coordinate upward edge from its lower endpoint.
    A coordinate rectangle is a product of a horizontal coordinate set and a
    vertical coordinate set. The $2\times3$ and
    $3\times2$ coordinate-product predicates assert only that these two
    coordinate sets have cardinalities $2,3$ and $3,2$, respectively.
    Contiguous coordinate rectangles are products of contiguous coordinate
    intervals, and the contiguous $2\times3$ and $3\times2$ predicates name the
    rectangular blocks used in the normal square-lattice theorem. The
    displayed region $R$ is the union of a $2\times3$ contiguous rectangle
    with the $3\times2$ contiguous rectangle occupying its upper two rows and
    extending one column to the right. The displayed region $S$ is the union
    of two $2\times3$ contiguous rectangles shifted by one horizontal
    coordinate, hence the full $3\times3$ square spanned by them. The
    displayed local-window model for $T$ is the complement, inside a local
    $5\times6$ coordinate window, of the vertical $2\times3$ edge block
    and the horizontal $3\times2$ edge block shown in the source picture.
    The finite-lattice complementary block is the whole finite square lattice
    with those two edge blocks removed; this is the block denoted $A_3$ in
    the proof of \cite[Theorem~3]{Molnar2018NormalPEPS}.
    The source shapes are:
    \begin{tenkzequation}
        % Source/formula verdict: Molnar et al., arXiv:1804.04964,
        % paper_normal.tex:1405--1429 and source PDF p. 15 depict
        % R=(2-4,2-4)\{(4,4)} and S=(2-4,2-4).
        % Ink-to-index verdict: every dot is an ordinary tensor; every lattice
        % edge is a virtual index, and the selected fill records membership only.
        % Structural count verdict (not mechanically emitted as a lattice
        % signature): each 5-by-5 panel has 25 sites, 40 internal virtual
        % contractions, and 20 cut virtual bonds; (4,4) remains ordinary.
        % Structurally derived boundary/type verdict: no physical index is
        % exposed, so each panel has (V,P)=(20,0).
        \tenkzkernel
        \begin{tenkz}[rows={ket,ket,ket,ket,ket}, cols=5,
            west=open, east=open, north=open, south=open]
            \tnmark[form=enclosure, slot=selected, tint, label pos=nw]
                {(2,2) .. (4,4) - (4,4)}{$R$}
        \end{tenkz}
        $\qquad\text{and}\qquad$
        \begin{tenkz}[rows={ket,ket,ket,ket,ket}, cols=5,
            west=open, east=open, north=open, south=open]
            \tnmark[form=enclosure, slot=selected, tint, label pos=nw]
                {(2,2) .. (4,4)}{$S$}
        \end{tenkz}
    \end{tenkzequation}
    \begin{tenkzequation}
        % Formula verdict: in row-column coordinates the exact local formal
        % region is T=(1-6,1-5) minus (3-5,1-2) and (2-3,3-5); the two
        % removed blocks are respectively 3-by-2 and 2-by-3 cell rectangles.
        % Ink-to-index verdict: every dot is an ordinary tensor, every lattice
        % edge is a virtual contraction, and the selected fill records only
        % membership in T; it adds no contraction or tensor operation.
        % Contraction/boundary verdict: the 6-by-5 panel has 49 internal
        % virtual contractions and 22 cut virtual bonds, with no exposed
        % physical indices, hence boundary type (V,P)=(22,0).
        % Source verdict: arXiv:1804.04964, paper_normal.tex:1430--1445
        % depicts the same 5-by-6 window complement; this panel uses the exact
        % cell set of normalSquareRegionT rather than the source's rounded ink.
        \tenkzkernel
        \begin{tenkz}[rows={ket,ket,ket,ket,ket,ket}, cols=5,
            west=open, east=open, north=open, south=open]
            \tnmark[form=enclosure, slot=selected, tint, label pos=n]
                {(1,1) .. (6,5) - (3,1) .. (5,2) - (2,3) .. (3,5)}{$T$}
        \end{tenkz}
    \end{tenkzequation}
    % Maintainer note: injectivity of this finite-lattice complementary block
    % is not yet proved; this definition only names the region.
\end{definition}

\begin{lemma}[Basic identities for square-lattice coordinate regions]%
    \label{lem:peps_squareLatticeCoordinateRegion_identities}
    \lean{TNLean.PEPS.mem_squareLatticeRectangle}
    \lean{TNLean.PEPS.mem_squareLatticeCoordinateInterval}
    \lean{TNLean.PEPS.mem_squareLatticeContiguousRectangle}
    \lean{TNLean.PEPS.mem_normalSquareRegionR}
    \lean{TNLean.PEPS.mem_normalSquareRegionS}
    \lean{TNLean.PEPS.mem_normalSquareRegionTVerticalBlock}
    \lean{TNLean.PEPS.mem_normalSquareRegionTHorizontalBlock}
    \lean{TNLean.PEPS.mem_normalSquareRegionTHole}
    \lean{TNLean.PEPS.mem_normalSquareRegionT}
    \lean{TNLean.PEPS.mem_normalSquareEdgeComplementRegion}
    \lean{TNLean.PEPS.isTwoByThreeContiguousSquareLatticeRectangle_of_bounds}
    \lean{TNLean.PEPS.isThreeByTwoContiguousSquareLatticeRectangle_of_bounds}
    \lean{TNLean.PEPS.card_squareLatticeRectangle}
    \lean{TNLean.PEPS.squareLatticeFull_eq_univ}
    \leanok
    \uses{def:peps_squareLatticeCoordinateRegions}
    Membership in a coordinate rectangle is coordinatewise membership, the
    contiguous coordinate interval has the expected coordinate inequalities,
    membership in a contiguous coordinate rectangle is the conjunction of the
    four bounding inequalities, and membership in $R$ is the disjunction of
    membership in the $2\times3$ lower-left piece or in the shifted
    $3\times2$ upper piece. Membership in $S$ is the disjunction of
    membership in the left $2\times3$ piece or in the horizontally shifted
    $2\times3$ piece. Membership in the local-window $T$ region is
    membership in the local $5\times6$ window together with nonmembership in
    the union of the two displayed edge blocks. Bounded contiguous coordinate
    rectangles of sizes $2\times3$ and $3\times2$ satisfy, respectively,
    the predicates for
    contiguous $2\times3$ and contiguous $3\times2$ square-lattice
    rectangles. The cardinality of a coordinate rectangle is the
    product of the two coordinate cardinalities, and the full coordinate
    rectangle is the whole finite vertex set.
\end{lemma}
\begin{proof}\leanok
    The coordinate-rectangle statements follow from membership and cardinality
    identities for Cartesian products of finite coordinate sets. The
    contiguous-rectangle statement then unfolds the two coordinate intervals.
    The membership statements for $R$, $S$, and the two edge blocks removed
    from $T$ additionally unfold membership in a union of two such
    rectangles, while membership in $T$ also unfolds the set difference from
    the local $5\times6$ window. The shape predicates follow by using the
    displayed lower-left corners as witnesses.
\end{proof}

\begin{lemma}[The normal regions $R$ and $S$ are unions of contiguous
        rectangles]%
    \label{lem:peps_normalSquareRegionRS_rectangular_decomposition}
    \lean{TNLean.PEPS.normalSquareRegionR_rectangular_decomposition}
    \lean{TNLean.PEPS.normalSquareRegionS_rectangular_decomposition}
    \leanok
    \uses{def:peps_squareLatticeCoordinateRegions,
        lem:peps_squareLatticeCoordinateRegion_identities}
    If the displayed $3\times3$ window containing $R$ and $S$ lies in
    the finite square lattice, then $R$ is the union of a contiguous
    $2\times3$ rectangle and a contiguous $3\times2$ rectangle, while
    $S$ is the union of two contiguous $2\times3$ rectangles.
\end{lemma}
\begin{proof}\leanok
    This is the defining description of $R$ and $S$, together with the
    coordinate bounds which make the displayed rectangles contiguous
    $2\times3$ and $3\times2$ rectangles.
\end{proof}

\begin{lemma}[The edge blocks removed from $T$ are contiguous rectangles]%
    \label{lem:peps_normalSquareRegionTHole_rectangular_decomposition}
    \lean{TNLean.PEPS.normalSquareRegionTVerticalBlock_rectangular}
    \lean{TNLean.PEPS.normalSquareRegionTHorizontalBlock_rectangular}
    \lean{TNLean.PEPS.normalSquareRegionTHole_rectangular_decomposition}
    \leanok
    \uses{def:peps_squareLatticeCoordinateRegions,
        lem:peps_squareLatticeCoordinateRegion_identities}
    If the two edge blocks cut out of the displayed $T$-region lie in the
    finite square lattice, then they are respectively a contiguous
    $2\times3$ rectangle and a contiguous $3\times2$ rectangle.
\end{lemma}
\begin{proof}\leanok
    This is the defining description of the two removed edge blocks, together
    with the coordinate bounds which make them valid contiguous rectangles.
\end{proof}

\begin{lemma}[The local $T$-region and its removed edge blocks fill the window]%
    \label{lem:peps_normalSquareRegionT_window_complement}
    \lean{TNLean.PEPS.normalSquareRegionTHole_subset_window}
    \lean{TNLean.PEPS.normalSquareRegionTVerticalBlock_disjoint_horizontalBlock}
    \lean{TNLean.PEPS.normalSquareRegionT_disjoint_THole}
    \lean{TNLean.PEPS.normalSquareRegionT_disjoint_verticalBlock}
    \lean{TNLean.PEPS.normalSquareRegionT_disjoint_horizontalBlock}
    \lean{TNLean.PEPS.normalSquareRegionT_union_THole}
    \lean{TNLean.PEPS.normalSquareRegionT_union_verticalBlock_union_horizontalBlock}
    \leanok
    \uses{def:peps_squareLatticeCoordinateRegions,
        lem:peps_squareLatticeCoordinateRegion_identities}
    The two edge blocks removed from the displayed local $T$-region lie inside
    the local $5\times6$ coordinate window. They are disjoint from each other
    and from this local $T$ region, and their union with it is exactly that
    local window.
\end{lemma}
\begin{proof}\leanok
    This follows by unfolding the definition of $T$ as the set-theoretic
    complement of the two displayed edge blocks inside the local window.
\end{proof}

\begin{lemma}[The edge-complementary block fills the finite lattice]%
    \label{lem:peps_normalSquareEdgeComplementRegion}
    \lean{TNLean.PEPS.normalSquareRegionT_subset_edgeComplementRegion}
    \lean{TNLean.PEPS.normalSquareEdgeComplementRegion_disjoint_THole}
    \lean{TNLean.PEPS.normalSquareEdgeComplementRegion_union_THole}
    \lean{TNLean.PEPS.normalSquareEdgeComplementRegion_union_verticalBlock_union_horizontalBlock}
    \leanok
    \uses{def:peps_squareLatticeCoordinateRegions,
        lem:peps_squareLatticeCoordinateRegion_identities}
    The finite-lattice complementary block around the edge is disjoint from
    the two displayed edge blocks, and its union with them is the full finite
    square lattice. The local-window $T$ region is contained in this
    finite-lattice complementary block.
\end{lemma}
\begin{proof}\leanok
    This is the set-theoretic complement of the two displayed edge blocks in
    the full finite square lattice. The containment of the local-window
    $T$ region follows because its points also avoid the two removed edge
    blocks.
\end{proof}

\begin{definition}[Rectangular covers of square-lattice regions]%
    \label{def:peps_normalSquareRectangleCover}
    \lean{TNLean.PEPS.SquareLatticeRectangleCover}
    \lean{TNLean.PEPS.NormalSquareRegionTRectangleCover}
    \lean{TNLean.PEPS.NormalSquareEdgeComplementRectangleCover}
    \lean{TNLean.PEPS.normalSquareVerticalEdgeComplementRegion}
    \lean{TNLean.PEPS.NormalSquareVerticalEdgeComplementRectangleCover}
    \leanok
    \uses{def:peps_squareLatticeCoordinateRegions,
        lem:peps_normalSquareEdgeComplementRegion}
    A rectangular cover of a square-lattice region is a nonempty finite family
    of regions whose union is exactly the given region, and each member is a
    contiguous $2\times3$ or $3\times2$ rectangle. The two cases used in the
    proof would be a cover of the local displayed $T$-region and a cover of the
    finite-lattice complementary block $A_3$. The panel below shows only the
    two permitted shapes of a cover member; it does not assert that these
    particular rectangles cover either region.
    \begin{tenkzequation}
        % Formula verdict: the selected cells form one 2-by-3 rectangle and
        % the secondary cells form one 3-by-2 rectangle.  This is a schematic
        % of the two allowed member shapes, not a union equation or a cover.
        % Ink-to-index verdict: every dot is an ordinary tensor and every
        % lattice edge is a virtual contraction; the outlined fills record
        % rectangle membership only and remain independently customizable.
        % Contraction/boundary verdict: the 4-by-4 panel has 24 internal
        % virtual contractions and 16 cut virtual bonds, with no exposed
        % physical indices, hence boundary type (V,P)=(16,0).
        % Source verdict: arXiv:1804.04964, paper_normal.tex:1405--1445 uses
        % 2-by-3 and 3-by-2 injective pieces.  The following formal lemma proves
        % the local no-cover obstruction for the present origin-based T and A3,
        % so no stronger geometric claim is encoded by this schematic.
        \tenkzkernel
        \begin{tenkz}[rows={ket,ket,ket,ket}, cols=4,
            west=open, east=open, north=open, south=open]
            \tnmark[form=enclosure, slot=selected]{(1,1) .. (2,3)}{$2\times3$}
            \tnmark[form=enclosure, slot=secondary]{(2,3) .. (4,4)}{$3\times2$}
        \end{tenkz}
    \end{tenkzequation}
\end{definition}
% Maintainer note: this definition records sufficient coordinate criteria.
% The concrete local T cover remains to be constructed.

\begin{lemma}[The present local $T$-window and $A_3$ cover obstructions]%
    \label{lem:peps_normalSquareRegionT_no_five_by_six_cover}
    \lean{TNLean.PEPS.not_squareLatticeRectangleCover_of_forced_points}
    \lean{TNLean.PEPS.not_normalSquareRegionT_rectangleCover_at_origin}
    \lean{TNLean.PEPS.not_normalSquareRegionT_rectangleCover_five_by_six}
    \lean{TNLean.PEPS.not_normalSquareVerticalRegionT_rectangleCover_at_origin}
    \lean{TNLean.PEPS.not_normalSquareVerticalRegionT_rectangleCover_six_by_five}
    \lean{TNLean.PEPS.not_normalSquareEdgeComplementRectangleCover_at_origin}
    \lean{TNLean.PEPS.not_normalSquareEdgeComplementRectangleCover_five_by_seven}
    \lean{TNLean.PEPS.not_normalSquareVerticalEdgeComplementRectangleCover_at_origin}
    \lean{TNLean.PEPS.not_normalSquareVerticalEdgeComplementRectangleCover_seven_by_five}
    \leanok
    \uses{def:peps_normalSquareRectangleCover,
        lem:peps_normalSquareRegionT_window_complement}
    The present origin-based local-window model for the displayed $T$-region
    admits no nonempty finite cover by contained contiguous $2\times3$ and
    $3\times2$ rectangles whenever the ambient lattice contains the
    $5\times6$ window. The origin-based rotated local-window model used for
    vertical edges has the analogous obstruction whenever the ambient lattice
    contains the $6\times5$ window. The normalized $5\times6$ and
    $6\times5$ cases are recorded as corollaries. The present origin-based
    horizontal edge-complement model has the analogous obstruction whenever
    the ambient lattice contains the $5\times7$ frame; the normalized
    $5\times7$ case is again recorded as a corollary. The present
    origin-based vertical edge-complement model has the corresponding
    obstruction whenever the ambient lattice contains the $7\times5$ frame;
    the normalized $7\times5$ case is recorded as a corollary.
\end{lemma}
% Maintainer note: this is a local-model obstruction, not a statement of the
% source theorem. It records that the current coordinate model is not yet the
% source cover needed for the $T$-injectivity sentence.
\begin{proof}\leanok
    The lower-left point of the local window belongs to the present $T$
    model. Any contiguous $2\times3$ rectangle containing it also contains
    a point of the removed vertical edge block, while any contiguous
    $3\times2$ rectangle containing it also contains a point of the removed
    horizontal edge block. The rotated statement is the same coordinate
    argument with the two shapes interchanged. The horizontal
    edge-complement statement uses the same lower-left point, and so does the
    normalized vertical edge-complement statement. These cases are all
    instances of the same point-forcing obstruction to an exact rectangular
    cover.
\end{proof}

\begin{lemma}[Injectivity from a rectangular cover]%
    \label{lem:peps_normalSquareRectangleCover_injective}
    \lean{TNLean.PEPS.NormalSquareLatticeRectangleInjectivityHypotheses.injective_of_rectangleCover}
    \leanok
    \uses{def:peps_normalSquareRectangleCover,
        def:peps_normalRectangleInjectivityHypotheses,
        lem:peps_regionInjectivityUnionClosure_finite}
    If a square-lattice region has a nonempty finite cover by contiguous
    $2\times3$ and $3\times2$ rectangles, then rectangular injectivity and
    finite union closure imply that the region is injective.
\end{lemma}
\begin{proof}\leanok
    Apply rectangular injectivity to every rectangle in the cover, then apply
    finite union closure to the nonempty family.
\end{proof}

\begin{lemma}[The normal regions $R$ and $S$ differ in one site]%
    \label{lem:peps_normalSquareRegionR_regionS_difference}
    \lean{TNLean.PEPS.normalSquareRegionR_subset_regionS}
    \lean{TNLean.PEPS.mem_normalSquareRegionS_sdiff_regionR}
    \leanok
    \uses{def:peps_squareLatticeCoordinateRegions,
        lem:peps_squareLatticeCoordinateRegion_identities}
    For the displayed regions $R$ and $S$ with the same lower-left corner,
    $R$ is contained in $S$. The difference $S\setminus R$ is exactly
    the lower-right site of the $3\times3$ square: in coordinates, the site
    one obtains by shifting two steps horizontally and zero steps vertically
    from the lower-left corner.
\end{lemma}
\begin{proof}\leanok
    Expand the two displayed unions. The only point of the second
    $2\times3$ piece of $S$ which is not already in the union defining
    $R$ is its lower-right point.
\end{proof}
